%----------------------------------------------------------------------------------------
\chapter{Introduzione}\label{chap:introduction}
%----------------------------------------------------------------------------------------
Il progetto di tesi è volto alla progettazione ex-novo del programma di vendita dell'azienda Cepi S.p.a. ,
azienda con una produzione di tipo \ac{ETO} specializzato nella progettazione di impianti di stoccaggio custom.
La tipologia di produzione \ac{ETO} ha portato ad un profondo studio della situazione odierna, delle sue esigenge attuali
e future tenendo aperte le possibilità di ampliamenti.
L'analisi ha comportato un sostanziale lavoro di programmazione e racccolta dati per coprire nel modo più
ampio i bisogni aziendali.
Questa prima fase è stata condotta inizialmente attraverso la selezione di figure chiave scelte per mansione o/e ruolo, che potessero
dare una visione dei bisogni del proprio o altrui reparto e successivamente sottoponendoli ad interviste.
Sono state condotte 14 interviste, di cui 8 in modalità vis a vis, mentre le restanti 6 attraverso la dimostrazione
dell'attuale applicativo utilizzato.
Il risultato delle interviste condotte, ha fatto emergere attuali e nuove necessità architteturali da affrontare
come il cross-platform, l'utilizzo offline del database e la necessità di costrurire offerte in modo più agile e strutturato.
Questa prima fase ha inoltre sottolineato come il progetto non debba solo sostituire l'attuale sistema,
ma debba inserirsi in un contesto di rinnovamento più ampio, coordinadosi con altri programmi aziendali e relativi database in ottica ``\ac{SRP}''
eleggendo dei master data che siano i soli resposabili della creazione di certi dati.
Il \ac{SRP} ha plasmato notevolmente la costruzione del database ricreando, dove possibile,
le strutture pre esistenti da cui si necessitava recuperare i dati mediante reverse engineering.~\cref{section:importdata}
L'intero sistema è stato pensato per appoggiarsi all'infrastruttura esistente aziendale, come SQLServer\cite{SQLServer},
ma studiando delle logiche di flessibilità e ampliamento dinamico studiando accuratamente le relazioni tra le tabelle.
Si è reso necessario cercare e scegliere inoltre un framework per la parte grafica che potesse soddisfare il cross-platform

essendo esso nelle richieste avanzate.
La costruzione delle interfacce grafiche, è stata supportata dal sito ``Figma''\cite{Figma} che ha permesso una progettazione %vedere se citarlo% 
in ottica mobile first.\ref{chap:View}
La scelta accurata dei framework e le strutture dati nella fase iniziale, ha pesato molto in termini di tempo e progettazione,
per garantire al richiedente un sistema che non rischiasse ne il debito tecnico, ne la preclusione a sviluppi futuri.
La pianificazione dei primi test con gli utenti, dei rilasci delle demo e del metodo di progettazione sono
state cruciali per dettare tempi e priorità degli item dovendo coprire in modo esaustivo almeno il core del progetto.
